%%%%%%%%%%%%%%%%%%%%%%%%%%%%%%%%%%%%%%%%%%%%%%%%%%%%%%%%%%%%%%%%%%%%%%%%%%%%%%%%
\section{Dealing with temporary Files}
{   
    \usebackgroundtemplate{
        \vbox to \paperheight{\vspace*{-1cm}\hbox to \paperwidth{\hfil\includegraphics[width=0.55\textwidth]{misc/broom_drawing.jpg}}\vfil}
    }
    \frame{
        \frametitle{Autodeleting temporary Output}
        \begin{mdframed}[tikzsetting={draw=white,fill=white,fill opacity=0.8,
                line width=0pt},backgroundcolor=none,leftmargin=0,
            rightmargin=150,innertopmargin=4pt,roundcorner=10pt]
            \tableofcontents[currentsection,sections={1-4},hideothersubsections]
            
        \end{mdframed}
        \vspace{1em}
        \tiny {water colour - C. Meesters - CC-BY-4.0}
    }
}


%%%%%%%%%%%%%%%%%%%%%%%%%%%%%%%%%%%%%%%%%%%%%%%%%%%%%%%%%%%%%%%%%%%%%%%%%%%%%%%%
\begin{frame}[fragile]
    \frametitle{Temporary Files}
    Let's face it: Most output files of a workflow, particularly more complex ones(!), do not need to be archived. They will only take up storage until someone takes the \includegraphics[height=\baselineskip]{misc/broom.png}.\newline
    Luckily, \Snakemake{} provides a way to mark a file as ``temporary'' and will delete it, when its not needed any longer.
    \pause
    \begin{question}
        The output of which rules in our workflow should be marked ``temporary''?
    \end{question}
\end{frame}

%%%%%%%%%%%%%%%%%%%%%%%%%%%%%%%%%%%%%%%%%%%%%%%%%%%%%%%%%%%%%%%%%%%%%%%%%%%%%%%%
\begin{frame}[fragile]
    \frametitle{Temporary Files - Reasoning}
    \begin{center}
        \begin{tabular}{|p{0.2\textwidth}|p{0.1\textwidth}|p{0.65\textwidth}}
            Rule & Temp? & Reason\\\hline
            \altverb{all} & \pause{\color{BrickRed}\bf No} & This rule desingates the final target(s).\\\hline
            \altverb{bwa_map} & \pause{\color{PineGreen}\bf Yes} & Mappings are derived from the intial input and can be created reproducibly again.\\\hline
            \altverb{samtools_sort} & \pause{\color{PineGreen}\bf Yes} & Same reason as with \altverb{bwa_map}\\\hline
            \altverb{samtools_index} & \pause{\color{PineGreen}\bf Yes} & Same reason as with \altverb{bwa_map}\\\hline
            \altverb{bcftools_call} & \pause{\color{BrickRed}\bf No} & This is a result file!\\\hline
            plotting scripts & \pause{\color{BrickRed}\bf No} & These produce result files! See \altverb{all}.\\\hline 
        \end{tabular}
    \end{center}
\end{frame} 

%%%%%%%%%%%%%%%%%%%%%%%%%%%%%%%%%%%%%%%%%%%%%%%%%%%%%%%%%%%%%%%%%%%%%%%%%%%%%%%%
\begin{frame}[fragile]
    \frametitle{\HandsOn{Temporary Files - Put to Work}}
    \Snakemake{}s way to mark outputs as temporary or intermediate is the \altverb{temp()}-function, e.\,g.:
    \begin{lstlisting}[language=Python,style=Python]
rule bwa_map:
    input:
        "data/genome.fa",
        get_bwa_map_input_fastqs
    output:
        @temp(@"mapped_reads/{sample}.bam"@)@
    ...
    \end{lstlisting}
    \pause
    Implement the \altverb{temp()} function for \altverb{bwa_map}`s output and run the workflow. Please run:
    \begin{lstlisting}[language=Bash, style=Shell]
$ du -sh mapped_reads # measures the size of the folder
$ snakemake -F # after adding temp()
$ du -sh mapped_reads # measures the size of the folder
    \end{lstlisting}
    \pause
    \bcquestion Do you see the benefit of temporary/intermediate files for your data management?
\end{frame}

%%%%%%%%%%%%%%%%%%%%%%%%%%%%%%%%%%%%%%%%%%%%%%%%%%%%%%%%%%%%%%%%%%%%%%%%%%%%%%%%
\begin{frame}[fragile]
    \frametitel{There is an Issue with my Data! $\ldots$}
    $\ldots$ how can I debug when \altverb{temp} is on?
    \pause
    \vfill
    \Snakemake offers a \altverb{--notemp} flag for that purpose.
    \vfill
    \pause
    \begin{alertblock}{Reproducibility and faulty Data}
        Reproducibilty does \textbf{not} mean, you data cannot be flawed. When you tinker with it, be sure to take note and be transparent about it!
    \end{alertblock}
\end{frame}


%%%%%%%%%%%%%%%%%%%%%%%%%%%%%%%%%%%%%%%%%%%%%%%%%%%%%%%%%%%%%%%%%%%%%%%%%%%%%%%%
\begin{frame}[fragile]
    \frametitle{The \altverb{protected()} Function}
    \begin{hint}
        This is just to demonstrate the complement to \altverb{temp()}. Sometimes you might want to protect files from accidental deletion. This can be done with the \altverb{protected()} from \Snakemake.
    \end{hint}
    \pause
    Files labeled "protected" are protected against accidental deletion and overwriting, by withdrawing the right to write. The implementation is as easy as:
    \begin{lstlisting}[language=Python,style=Python]
rule plot_quals:
    ...
    output: 
        protected("calls/quals.png")
...
    \end{lstlisting}
    \pause
    \begin{hint}
        We will \emph{NOT} do this, as it will interfere with our repeated executions.
    \end{hint}	
\end{frame} 